% !TeX root = part01.tex

\documentclass[../final.tex]{subfiles}

\begin{document}

% ============================================================
% Практическое занятие 2
% ============================================================

\practice{2}{08.09.2026}{Классические и квантовые модели атома}

% ============================================================
% Задача 3. Классическое время падения электрона на ядро
% ============================================================

\section{Задача 3. Классическое время падения электрона на ядро}


% ------------------------------------------------------------
% Условие
% ------------------------------------------------------------

\rtxt{3}{\textbf{Условие:}}
Оценить время, за которое электрон, движущийся вокруг ядра водорода
по орбите радиусом
%
\begin{equation*}
    r_0=0.5\,\text{\AA},
\end{equation*}
%
упал бы на ядро, если бы он терял энергию на излучение
в соответствии с классической теорией.


% ------------------------------------------------------------
% Схема
% ------------------------------------------------------------

\begin{rbox}[left=18pt,right=18pt,top=14pt,bottom=10pt]{[]}
    \centering

    \begin{tikzpicture}[
        >=Latex,
        every node/.style={text=rtext}
    ]
        \coordinate (O) at (0,0);

        % Орбита
        \draw[
            draw=rtext,
            line width=0.9pt
        ]
        (O) circle (1.45);

        % Ядро
        \fill[rtext] (O) circle (2.4pt);
        \node[below=5pt] at (O) {$p^+$};

        % Электрон
        \coordinate (E) at (32:1.45);
        \fill[rc3] (E) circle (2.6pt);
        \node[right=5pt] at (E) {$e^-$};

        % Радиус
        \draw[
            draw=rtext,
            line width=0.8pt
        ]
        (O) -- (E);

        \node[above left] at (0.62,0.42) {$r$};

        % Направление движения
        \draw[
            ->,
            draw=rtext,
            line width=0.9pt
        ]
        (58:1.45)
        arc[
            start angle=58,
            end angle=118,
            radius=1.45
        ];
    \end{tikzpicture}

    \captionsetup{
        font=footnotesize,
        labelfont=bf,
        skip=5pt
    }

    \captionof{figure}{
        Классическая планетарная модель атома водорода.
        Электрон движется по круговой орбите радиуса $r$
        вокруг ядра.
    }
    \label{fig:classical_hydrogen_fall}
\end{rbox}


% ------------------------------------------------------------
% Энергия электрона
% ------------------------------------------------------------

Полная энергия электрона на круговой орбите равна
%
\begin{equation*}
    E
    =
    \frac{mv^2}{2}
    -
    \frac{e^2}{r}.
\end{equation*}

Кулоновская сила является центростремительной:
%
\begin{equation*}
    \frac{mv^2}{r}
    =
    \frac{e^2}{r^2}.
\end{equation*}

Следовательно,
%
\begin{equation*}
    v^2
    =
    \frac{e^2}{mr}.
\end{equation*}

Подставляя это выражение в полную энергию, получаем
%
\begin{equation*}
    E
    =
    \frac{m}{2}\frac{e^2}{mr}
    -
    \frac{e^2}{r}
    =
    \frac{e^2}{2r}
    -
    \frac{e^2}{r}
    =
    -\frac{e^2}{2r}.
\end{equation*}


% ------------------------------------------------------------
% Мощность излучения
% ------------------------------------------------------------

Согласно классической электродинамике, ускоренно движущийся заряд
излучает. Мощность излучения определяется формулой Лармора:
%
\begin{equation*}
    P
    =
    \frac{2e^2}{3c^3}a^2.
\end{equation*}

При движении по окружности
%
\begin{equation*}
    a
    =
    \frac{v^2}{r}.
\end{equation*}

Поскольку
%
\begin{equation*}
    v^2
    =
    \frac{e^2}{mr},
\end{equation*}
%
то
%
\begin{equation*}
    a
    =
    \frac{e^2}{mr^2}.
\end{equation*}

Следовательно,
%
\begin{equation*}
    P
    =
    \frac{2e^2}{3c^3}
    \left(
        \frac{e^2}{mr^2}
    \right)^2
    =
    \frac{2e^6}{3m^2c^3r^4}.
\end{equation*}


% ------------------------------------------------------------
% Уменьшение радиуса орбиты
% ------------------------------------------------------------

Энергия электрона уменьшается вследствие излучения, поэтому
%
\begin{equation*}
    P
    =
    -\frac{dE}{dt}.
\end{equation*}

Так как
%
\begin{equation*}
    E
    =
    -\frac{e^2}{2r},
\end{equation*}
%
имеем
%
\begin{equation*}
    \frac{dE}{dt}
    =
    \frac{e^2}{2r^2}\frac{dr}{dt}.
\end{equation*}

Тогда
%
\begin{equation*}
    \frac{e^2}{2r^2}\frac{dr}{dt}
    =
    -\frac{2e^6}{3m^2c^3r^4}.
\end{equation*}

Отсюда
%
\begin{equation*}
    r^2\frac{dr}{dt}
    =
    -\frac{4e^4}{3m^2c^3}.
\end{equation*}

Интегрируем от начального радиуса $r_0$ до $r=0$
и от $t=0$ до времени падения $\tau$:
%
\begin{equation*}
    \int_{r_0}^{0}r^2\,dr
    =
    -\frac{4e^4}{3m^2c^3}
    \int_0^\tau dt.
\end{equation*}

Получаем
%
\begin{equation*}
    -\frac{r_0^3}{3}
    =
    -\frac{4e^4}{3m^2c^3}\tau.
\end{equation*}

Следовательно,

\[
    \boxed{
        \tau
        =
        \frac{m^2c^3r_0^3}{4e^4}
    }.
\]

Для
%
\begin{equation*}
    r_0
    =
    0.5\,\text{\AA}
    =
    0.5\cdot10^{-8}\,\text{см}
\end{equation*}
%
получаем

\[
    \boxed{
        \tau
        \approx
        1.3\cdot10^{-11}\,\text{с}
    }.
\]

Таким образом, в рамках классической электродинамики атом водорода
должен был бы разрушиться за чрезвычайно малое время.


% ============================================================
% Задача 4. Квантование движения в поле U = chi r^2 / 2
% ============================================================

\section{Задача 4. Квантование движения в поле
\texorpdfstring{$U=\frac12\chi r^2$}{U = chi r²/2}}


% ------------------------------------------------------------
% Условие
% ------------------------------------------------------------

\rtxt{3}{\textbf{Условие:}}
Частица массой $m$ движется по круговой орбите в поле
%
\begin{equation*}
    U(r)
    =
    \frac{1}{2}\chi r^2.
\end{equation*}
%
Найти с помощью боровского условия квантования разрешённые
уровни энергии и соответствующие радиусы орбит.


% ------------------------------------------------------------
% Уравнение движения
% ------------------------------------------------------------

Сила определяется через потенциальную энергию:
%
\begin{equation*}
    F
    =
    -\frac{dU}{dr}
    =
    -\chi r.
\end{equation*}

Для круговой орбиты модуль центростремительной силы равен
%
\begin{equation*}
    \frac{mv^2}{r}
    =
    \chi r.
\end{equation*}

Следовательно,
%
\begin{equation*}
    mv^2
    =
    \chi r^2,
    \qquad
    v
    =
    r\sqrt{\frac{\chi}{m}}.
\end{equation*}


% ------------------------------------------------------------
% Условие Бора
% ------------------------------------------------------------

Используем условие квантования момента импульса:
%
\begin{equation*}
    L
    =
    mvr
    =
    n\hbar,
    \qquad
    n=1,2,3,\ldots
\end{equation*}

Подставляя скорость,
%
\begin{equation*}
    mr^2
    \sqrt{\frac{\chi}{m}}
    =
    n\hbar.
\end{equation*}

Отсюда
%
\begin{equation*}
    r^2
    =
    \frac{n\hbar}{\sqrt{m\chi}}.
\end{equation*}

Следовательно, разрешённые радиусы орбит равны

\[
    \boxed{
        r_n
        =
        \left(
            \frac{n^2\hbar^2}{m\chi}
        \right)^{1/4}
    }.
\]


% ------------------------------------------------------------
% Энергия
% ------------------------------------------------------------

Полная энергия частицы
%
\begin{equation*}
    E
    =
    \frac{mv^2}{2}
    +
    \frac{\chi r^2}{2}.
\end{equation*}

Из уравнения движения
%
\begin{equation*}
    mv^2
    =
    \chi r^2,
\end{equation*}
%
поэтому
%
\begin{equation*}
    E
    =
    \frac{\chi r^2}{2}
    +
    \frac{\chi r^2}{2}
    =
    \chi r^2.
\end{equation*}

Подставляя $r_n$,
%
\begin{equation*}
    E_n
    =
    \chi
    \frac{n\hbar}{\sqrt{m\chi}}
    =
    n\hbar
    \sqrt{\frac{\chi}{m}}.
\end{equation*}

Введём собственную частоту осциллятора
%
\begin{equation*}
    \omega
    =
    \sqrt{\frac{\chi}{m}}.
\end{equation*}

Тогда

\[
    \boxed{
        E_n
        =
        n\hbar\omega
    }.
\]

Для сравнения, точное решение уравнения Шрёдингера для гармонического
осциллятора даёт
%
\begin{equation*}
    E_n^{\text{QM}}
    =
    \hbar\omega
    \left(
        n+\frac12
    \right),
    \qquad
    n=0,1,2,\ldots
\end{equation*}

Таким образом, боровское квантование правильно воспроизводит
расстояние между соседними уровнями
%
\begin{equation*}
    \Delta E
    =
    \hbar\omega,
\end{equation*}
%
но не даёт нулевую энергию $\hbar\omega/2$.


% ============================================================
% Задача 5. Спектр атома водорода
% ============================================================

\section{Задача 5. Спектр атома водорода}


% ------------------------------------------------------------
% Условие
% ------------------------------------------------------------

\rtxt{3}{\textbf{Условие:}}
Пренебрегая спин-орбитальным взаимодействием для атомарного водорода,
вычислить:

\begin{enumerate}
    \item[а)]
    в каких пределах должна лежать энергия бомбардирующих электронов,
    чтобы спектр излучения атома имел только три линии;
    указать их длины волн;

    \item[б)]
    минимальную разрешающую способность спектрометра
    %
    \begin{equation*}
        R
        =
        \frac{\lambda}{\delta\lambda},
    \end{equation*}
    %
    при которой можно разрешить первые $20$ линий серии Бальмера.
\end{enumerate}


% ------------------------------------------------------------
% Часть а
% ------------------------------------------------------------

\subsection*{а) Три линии в спектре}

Энергетические уровни атома водорода определяются формулой
%
\begin{equation*}
    E_n
    =
    -\frac{Ry}{n^2}.
\end{equation*}

Для первых четырёх уровней:
%
\begin{equation*}
    E_1=-13.6\,\text{эВ},
    \qquad
    E_2=-3.4\,\text{эВ},
    \qquad
    E_3\approx-1.51\,\text{эВ},
    \qquad
    E_4=-0.85\,\text{эВ}.
\end{equation*}


% ------------------------------------------------------------
% Схема переходов
% ------------------------------------------------------------

\begin{rbox}[left=18pt,right=18pt,top=14pt,bottom=10pt]{[]}
    \centering

    \begin{tikzpicture}[
        >=Latex,
        every node/.style={text=rtext}
    ]
        % Уровни
        \draw[
            draw=rtext,
            line width=0.9pt
        ]
        (-2.5,0) -- (2.5,0);

        \draw[
            draw=rtext,
            line width=0.9pt
        ]
        (-2.5,1.5) -- (2.5,1.5);

        \draw[
            draw=rtext,
            line width=0.9pt
        ]
        (-2.5,2.4) -- (2.5,2.4);

        % Подписи уровней
        \node[right] at (2.5,0)
            {$n=1,\quad E_1=-13.6\,\text{эВ}$};

        \node[right] at (2.5,1.5)
            {$n=2,\quad E_2=-3.4\,\text{эВ}$};

        \node[right] at (2.5,2.4)
            {$n=3,\quad E_3\approx-1.51\,\text{эВ}$};

        % Переход 3 -> 1
        \draw[
            ->,
            draw=rc3,
            line width=1.0pt
        ]
        (-1.25,2.33) -- (-1.25,0.07);

        \node[left] at (-1.25,1.15)
            {$3\rightarrow1$};

        % Переход 3 -> 2
        \draw[
            ->,
            draw=rc3,
            line width=1.0pt
        ]
        (0,2.33) -- (0,1.57);

        \node[right] at (0,1.95)
            {$3\rightarrow2$};

        % Переход 2 -> 1
        \draw[
            ->,
            draw=rc3,
            line width=1.0pt
        ]
        (1.25,1.43) -- (1.25,0.07);

        \node[right] at (1.25,0.75)
            {$2\rightarrow1$};
    \end{tikzpicture}

    \captionsetup{
        font=footnotesize,
        labelfont=bf,
        skip=5pt
    }

    \captionof{figure}{
        Три возможных спектральных перехода при возбуждении
        атома водорода не выше уровня $n=3$.
    }
    \label{fig:hydrogen_three_lines}
\end{rbox}


Чтобы были доступны состояния вплоть до $n=3$, но уровень $n=4$
ещё не возбуждался, необходимо
%
\begin{equation*}
    \Delta E_{1\rightarrow3}
    \leq
    E
    <
    \Delta E_{1\rightarrow4}.
\end{equation*}

Для перехода $1\rightarrow3$
%
\begin{equation*}
    \Delta E_{1\rightarrow3}
    =
    E_3-E_1
    =
    Ry
    \left(
        1-\frac{1}{9}
    \right)
    =
    \frac{8}{9}Ry.
\end{equation*}

Для перехода $1\rightarrow4$
%
\begin{equation*}
    \Delta E_{1\rightarrow4}
    =
    E_4-E_1
    =
    Ry
    \left(
        1-\frac{1}{16}
    \right)
    =
    \frac{15}{16}Ry.
\end{equation*}

Следовательно,

\[
    \boxed{
        \frac{8}{9}Ry
        \leq
        E
        <
        \frac{15}{16}Ry
    }.
\]

При $Ry=13.6\,\text{эВ}$:
%
\begin{equation*}
    \boxed{
        12.09\,\text{эВ}
        \lesssim
        E
        <
        12.75\,\text{эВ}
    }.
\end{equation*}


% ------------------------------------------------------------
% Длины волн
% ------------------------------------------------------------

Энергия фотона при переходе $n_i\rightarrow n_f$ равна
%
\begin{equation*}
    h\nu
    =
    Ry
    \left(
        \frac{1}{n_f^2}
        -
        \frac{1}{n_i^2}
    \right),
    \qquad
    \nu
    =
    \frac{c}{\lambda}.
\end{equation*}

Поэтому
%
\begin{equation*}
    \lambda
    =
    \frac{hc}{
        Ry
        \left(
            \dfrac{1}{n_f^2}
            -
            \dfrac{1}{n_i^2}
        \right)
    }.
\end{equation*}

Для перехода $3\rightarrow1$:
%
\begin{equation*}
    \boxed{
        \lambda_{3\rightarrow1}
        \approx
        102.6\,\text{нм}
    }.
\end{equation*}

Для перехода $2\rightarrow1$:
%
\begin{equation*}
    \boxed{
        \lambda_{2\rightarrow1}
        \approx
        121.6\,\text{нм}
    }.
\end{equation*}

Обе эти линии относятся к серии Лаймана.

Для перехода $3\rightarrow2$:
%
\begin{equation*}
    \boxed{
        \lambda_{3\rightarrow2}
        \approx
        656.5\,\text{нм}
    }.
\end{equation*}

Это головная линия серии Бальмера.


% ------------------------------------------------------------
% Часть б
% ------------------------------------------------------------

\subsection*{б) Разрешение первых 20 линий серии Бальмера}

Для серии Бальмера конечный уровень равен
%
\begin{equation*}
    n_f=2.
\end{equation*}

Первые двадцать линий соответствуют переходам
%
\begin{equation*}
    n_i
    =
    3,4,\ldots,22.
\end{equation*}

С ростом $n_i$ соседние линии сближаются. Поэтому наиболее трудно
разрешить две последние линии:
%
\begin{equation*}
    21\rightarrow2,
    \qquad
    22\rightarrow2.
\end{equation*}


% ------------------------------------------------------------
% Схема близких линий
% ------------------------------------------------------------

\begin{rbox}[left=18pt,right=18pt,top=14pt,bottom=10pt]{[]}
    \centering

    \begin{tikzpicture}[
        >=Latex,
        every node/.style={text=rtext}
    ]
        % n = 22
        \draw[
            draw=rc3,
            line width=1.0pt
        ]
        (-2.4,1.05) -- (2.0,1.05);

        \node[right] at (2.0,1.05)
            {$n=22$};

        % n = 21
        \draw[
            draw=rtext,
            line width=1.0pt
        ]
        (-2.4,0.55) -- (2.0,0.55);

        \node[right] at (2.0,0.55)
            {$n=21$};

        % Разность
        \draw[
            <->,
            draw=rtext,
            line width=0.8pt
        ]
        (-1.7,0.58) -- (-1.7,1.02);

        \node[left] at (-1.7,0.8)
            {$\delta E$};
    \end{tikzpicture}

    \captionsetup{
        font=footnotesize,
        labelfont=bf,
        skip=5pt
    }

    \captionof{figure}{
        Наиболее близко расположенные верхние уровни,
        определяющие минимальную разрешающую способность
        для первых двадцати линий серии Бальмера.
    }
    \label{fig:balmer_resolution_levels}
\end{rbox}


Для перехода $21\rightarrow2$:
%
\begin{equation*}
    \lambda_{21}
    =
    \frac{hc}{
        Ry
        \left(
            \dfrac14
            -
            \dfrac{1}{21^2}
        \right)
    }
    \approx
    368.05\,\text{нм}.
\end{equation*}

Для перехода $22\rightarrow2$:
%
\begin{equation*}
    \lambda_{22}
    =
    \frac{hc}{
        Ry
        \left(
            \dfrac14
            -
            \dfrac{1}{22^2}
        \right)
    }
    \approx
    367.75\,\text{нм}.
\end{equation*}

Разность длин волн
%
\begin{equation*}
    \delta\lambda
    =
    \lambda_{21}
    -
    \lambda_{22}
    \approx
    0.30\,\text{нм}.
\end{equation*}

Минимальная разрешающая способность спектрометра:
%
\begin{equation*}
    R
    =
    \frac{\lambda}{\delta\lambda}
    \approx
    \frac{368}{0.30}
    \approx
    1.23\cdot10^3.
\end{equation*}

Следовательно,

\[
    \boxed{
        R_{\min}
        \approx
        1.23\cdot10^3
        \approx
        1227
    }.
\]


% ============================================================
% Задача 6. Мезоатом водорода
% ============================================================

\section{Задача 6. Мезоатом водорода}


% ------------------------------------------------------------
% Условие
% ------------------------------------------------------------

\rtxt{3}{\textbf{Условие:}}
Вычислить для мезоатома водорода, если масса мезона составляет
$207$ масс электрона:

\begin{enumerate}
    \item[а)]
    эффективный радиус Бора;

    \item[б)]
    постоянную Ридберга и длину волны резонансной линии;

    \item[в)]
    энергии связи основных состояний, когда ядром является
    протон или дейтон; сравнить изотопический сдвиг
    со сдвигом в атоме водорода.
\end{enumerate}

Масса мюона:
%
\begin{equation*}
    m_\mu
    =
    207m_e.
\end{equation*}


% ------------------------------------------------------------
% Схема мезоатома
% ------------------------------------------------------------

\begin{rbox}[left=18pt,right=18pt,top=14pt,bottom=10pt]{[]}
    \centering

    \begin{tikzpicture}[
        >=Latex,
        every node/.style={text=rtext}
    ]
        \coordinate (O) at (0,0);

        % Орбита
        \draw[
            draw=rtext,
            line width=0.9pt
        ]
        (O) circle (1.45);

        % Ядро
        \fill[rtext] (O) circle (2.5pt);
        \node[below=5pt] at (O) {$M$};

        % Мюон
        \coordinate (MU) at (35:1.45);
        \fill[rc3] (MU) circle (2.7pt);
        \node[right=5pt] at (MU) {$\mu^-$};

        % Радиус
        \draw[
            draw=rtext,
            line width=0.8pt
        ]
        (O) -- (MU);

        \node[above left] at (0.58,0.43)
            {$a_B^\mu$};

        % Направление
        \draw[
            ->,
            draw=rtext,
            line width=0.9pt
        ]
        (62:1.45)
        arc[
            start angle=62,
            end angle=122,
            radius=1.45
        ];
    \end{tikzpicture}

    \captionsetup{
        font=footnotesize,
        labelfont=bf,
        skip=5pt
    }

    \captionof{figure}{
        Мезоатом водорода: вместо электрона вокруг положительно
        заряженного ядра находится отрицательный мюон.
    }
    \label{fig:muonic_hydrogen}
\end{rbox}


% ------------------------------------------------------------
% Приведённая масса
% ------------------------------------------------------------

Для системы из мюона и ядра массы $M$ необходимо использовать
приведённую массу
%
\begin{equation*}
    \mu
    =
    \frac{m_\mu M}{m_\mu+M}.
\end{equation*}

Именно приведённая масса должна использоваться вместо массы электрона
во всех формулах водородоподобного атома.


% ------------------------------------------------------------
% Часть а
% ------------------------------------------------------------

\subsection*{а) Эффективный радиус Бора}

Обычный боровский радиус равен
%
\begin{equation*}
    a_0
    =
    \frac{\hbar^2}{m_e e^2}.
\end{equation*}

Для мезоатома
%
\begin{equation*}
    a_B^\mu
    =
    \frac{\hbar^2}{\mu e^2}
    =
    a_0\frac{m_e}{\mu}.
\end{equation*}

Подставляя выражение для приведённой массы,
%
\begin{equation*}
    a_B^\mu
    =
    a_0m_e
    \frac{m_\mu+M}{m_\mu M}
    =
    a_0
    \frac{m_e}{m_\mu}
    \left(
        1+\frac{m_\mu}{M}
    \right).
\end{equation*}

Для протона
%
\begin{equation*}
    M=m_p\approx1836m_e.
\end{equation*}

Тогда
%
\begin{equation*}
    \frac{\mu_H}{m_e}
    =
    \frac{207\cdot1836}{207+1836}
    \approx
    186.
\end{equation*}

Следовательно,
%
\begin{equation*}
    a_B^{\mu H}
    =
    \frac{a_0}{186}
    \approx
    \frac{0.529\,\text{\AA}}{186}
    \approx
    2.84\cdot10^{-3}\,\text{\AA}.
\end{equation*}

Таким образом,

\[
    \boxed{
        a_B^{\mu H}
        \approx
        2.8\cdot10^{-3}\,\text{\AA}
    }.
\]

Мезоатом примерно в $186$ раз меньше обычного атома водорода.


% ------------------------------------------------------------
% Часть б
% ------------------------------------------------------------

\subsection*{б) Постоянная Ридберга и резонансная линия}

Ридберговская энергия пропорциональна приведённой массе:
%
\begin{equation*}
    Ry_\mu
    =
    Ry
    \frac{\mu}{m_e}.
\end{equation*}

Для мезоатома водорода
%
\begin{equation*}
    Ry_{\mu H}
    \approx
    13.6\,\text{эВ}\cdot186
    \approx
    2.53\cdot10^3\,\text{эВ}.
\end{equation*}

Следовательно,

\[
    \boxed{
        Ry_{\mu H}
        \approx
        2.53\,\text{кэВ}
    }.
\]

В спектроскопических единицах постоянная Ридберга увеличивается
во столько же раз:
%
\begin{equation*}
    R_{\mu H}
    \approx
    186R_\infty
    \approx
    2.04\cdot10^7\,\text{см}^{-1}.
\end{equation*}

Резонансная линия соответствует переходу
%
\begin{equation*}
    2\rightarrow1.
\end{equation*}

Её энергия равна
%
\begin{equation*}
    \Delta E_{2\rightarrow1}
    =
    Ry_{\mu H}
    \left(
        1-\frac14
    \right)
    =
    \frac34Ry_{\mu H}.
\end{equation*}

Численно
%
\begin{equation*}
    \Delta E_{2\rightarrow1}
    \approx
    \frac34
    \cdot
    2.53\,\text{кэВ}
    \approx
    1.90\,\text{кэВ}.
\end{equation*}

Длина волны
%
\begin{equation*}
    \lambda_{\mathrm{res}}
    =
    \frac{hc}{\Delta E_{2\rightarrow1}}
    \approx
    0.65\,\text{нм}.
\end{equation*}

Следовательно,

\[
    \boxed{
        \lambda_{\mathrm{res}}
        \approx
        0.65\,\text{нм}
        \approx
        6.5\,\text{\AA}
    }.
\]


% ------------------------------------------------------------
% Часть в
% ------------------------------------------------------------

\subsection*{в) Основные состояния и изотопический сдвиг}

Энергия основного состояния водородоподобной системы равна
%
\begin{equation*}
    E_1
    =
    -Ry_\mu.
\end{equation*}


% ------------------------------------------------------------
% Мюонный водород
% ------------------------------------------------------------

Для ядра-протона
%
\begin{equation*}
    \mu_H
    =
    \frac{m_\mu m_p}{m_\mu+m_p}
    \approx
    186.0m_e.
\end{equation*}

Поэтому
%
\begin{equation*}
    E_1^{\mu H}
    =
    -Ry
    \frac{\mu_H}{m_e}
    \approx
    -2.53\cdot10^3\,\text{эВ}.
\end{equation*}

То есть энергия связи основного состояния

\[
    \boxed{
        |E_1^{\mu H}|
        \approx
        2.53\,\text{кэВ}
    }.
\]


% ------------------------------------------------------------
% Мюонный дейтерий
% ------------------------------------------------------------

Для дейтрона
%
\begin{equation*}
    m_d
    \approx
    3670m_e.
\end{equation*}

Приведённая масса:
%
\begin{equation*}
    \mu_D
    =
    \frac{m_\mu m_d}{m_\mu+m_d}
    =
    \frac{207\cdot3670}{207+3670}m_e
    \approx
    196.0m_e.
\end{equation*}

Следовательно,
%
\begin{equation*}
    E_1^{\mu D}
    =
    -Ry
    \frac{\mu_D}{m_e}
    \approx
    -2.66\cdot10^3\,\text{эВ}.
\end{equation*}

То есть

\[
    \boxed{
        |E_1^{\mu D}|
        \approx
        2.66\,\text{кэВ}
    }.
\]


% ------------------------------------------------------------
% Изотопический сдвиг
% ------------------------------------------------------------

Изотопический сдвиг основных состояний мезоатомов:
%
\begin{equation*}
    \Delta E_\mu
    =
    |E_1^{\mu D}|
    -
    |E_1^{\mu H}|
    =
    Ry
    \frac{\mu_D-\mu_H}{m_e}.
\end{equation*}

Численно
%
\begin{equation*}
    \boxed{
        \Delta E_\mu
        \approx
        1.35\cdot10^2\,\text{эВ}
    }.
\end{equation*}

Для обычного атома водорода необходимо использовать приведённые массы
электрона с протоном и дейтроном:
%
\begin{equation*}
    \mu_H^{(e)}
    =
    \frac{m_em_p}{m_e+m_p},
    \qquad
    \mu_D^{(e)}
    =
    \frac{m_em_d}{m_e+m_d}.
\end{equation*}

Соответствующий изотопический сдвиг основных состояний равен
%
\begin{equation*}
    \Delta E_e
    =
    Ry
    \frac{
        \mu_D^{(e)}
        -
        \mu_H^{(e)}
    }{m_e}
    \approx
    3.7\cdot10^{-3}\,\text{эВ}.
\end{equation*}

Следовательно,
%
\begin{equation*}
    \frac{\Delta E_\mu}{\Delta E_e}
    \approx
    3.6\cdot10^4.
\end{equation*}

Таким образом,

\[
    \boxed{
        \Delta E_\mu
        \gg
        \Delta E_e
    }.
\]

Изотопический сдвиг в мезоатоме значительно больше, поскольку
масса мюона уже не является пренебрежимо малой по сравнению
с массой ядра.
\end{document}